\clearpage
\section{Конструкторская часть}

\subsection{Требования к программе}

К программе предъявляются следующие требования:
\begin{itemize}
    \item реализации алгоритмов принимают натуральные числа $s$ и $k$, \mbox{$k\le19$},
    и возвращают количество $k$-значных натуральных чисел с суммой цифр $s$;
    \item рекурсивная реализация не использует циклы, массивы, строки и срезы;
    \item программа выводит таблицу значений $F(i,j)$ для \mbox{$0\le i\le k$},
    \mbox{$0\le j\le s$} и результаты обоих алгоритмов.
\end{itemize}

\subsection{Нерекурсивный алгоритм}

Нерекурсивный алгоритм заполняет таблицу значений $F$ по формуле~\eqref{eq:f}:
строки~--- сверху вниз по возрастанию $i$, ячейки строки~--- слева направо
по возрастанию $j$. Алгоритм получает числа $k$ и $s$ и возвращает $F[k][s]$.
Псевдокод алгоритма приведён в листинге~\ref{lst:iter} «Псевдокод
нерекурсивного алгоритма», схема показана на рисунке~\ref{fig:scheme-iter}.

\bigskip
\noindent\begin{minipage}{\dimexpr\linewidth-0.5mm\relax}
\begin{lstlisting}[language=Python,caption={Псевдокод нерекурсивного алгоритма},label={lst:iter}]
def count(k, s):
    F = [[0] * (s + 1) for _ in range(k + 1)]
    F[0][0] = 1
    for i in range(1, k + 1):
        for j in range(1, s + 1):
            for c in range(0, min(9, j) + 1):
                F[i][j] += F[i - 1][j - c]
    return F[k][s]
\end{lstlisting}
\end{minipage}
\medskip

\begin{figure}[H]
    \centering
    \begin{tikzpicture}[node distance=5mm]
\node[terminator] (start) {Начало};
\node[process,below=of start] (b2) {Создать таблицу $F$ размером $(k+1)\times(s+1)$, заполненную нулями};
\node[process,below=of b2] (b3) {$F[0][0]=1$};
\node[loop start,below=of b3] (b4) {Цикл I\\$i$ от 1 до $k$};
\node[loop start,below=of b4] (b5) {Цикл J\\$j$ от 1 до $s$};
\node[loop start,below=of b5] (b6) {Цикл C\\$c$ от 0 до $\min(9,j)$};
\node[process,below=of b6] (b7) {$F[i][j]=F[i][j]+F[i-1][j-c]$};
\node[loop end,below=of b7] (e6) {Цикл C};
\node[loop end,below=of e6] (e5) {Цикл J};
\node[loop end,below=of e5] (e4) {Цикл I};
\node[process,below=of e4] (b8) {Вернуть $F[k][s]$};
\node[terminator,below=of b8] (end) {Конец};
\foreach \a/\b in {start/b2,b2/b3,b3/b4,b4/b5,b5/b6,b6/b7,b7/e6,e6/e5,e5/e4,e4/b8,b8/end}
    \draw[flow] (\a) -- (\b);
\end{tikzpicture}

    \caption{Схема нерекурсивного алгоритма}
    \label{fig:scheme-iter}
\end{figure}

\subsection{Рекурсивный алгоритм}

Рекурсивный алгоритм вычисляет $F(k,s)$ по формулам~\eqref{eq:f}
и~\eqref{eq:s} без сохранения промежуточных значений: перебор последней цифры
выполняется рекурсивной функцией $S$. Схемы функций $F$ и $S$ показаны
на рисунках~\ref{fig:scheme-f} и~\ref{fig:scheme-s}, псевдокод обеих функций
приведён в листинге~\ref{lst:rec} «Псевдокод рекурсивного алгоритма».

\begin{figure}[H]
    \centering
    \begin{tikzpicture}[x=1mm,y=-1mm]
\node[terminator] (start) at (0,0) {$F(i,j)$};
\node[sdecision] (b1) at (0,22) {$i=0$};
\node[sdecision] (b2) at (-40,44) {$j=0$};
\node[sdecision] (b5) at (40,44) {$j=0$};
\node[sprocess] (b3) at (-60,66) {$r=1$};
\node[sprocess] (b4) at (-20,66) {$r=0$};
\node[sprocess] (b6) at (20,66) {$r=0$};
\node[spredefined] (b7) at (60,66) {$r=S(i,j,0)$};
\node[terminator] (end) at (0,104) {Возврат: $r$};
\draw[flow] (start) -- (b1);
\draw[flow] (b1.west) -- node[lbl,above] {да} (-40,22) -- (b2.north);
\draw[flow] (b1.east) -- node[lbl,above] {нет} (40,22) -- (b5.north);
\draw[flow] (b2.west) -- (-60,44) -- node[lbl,left,pos=0.35] {да} (b3.north);
\draw[flow] (b2.east) -- (-20,44) -- node[lbl,right,pos=0.35] {нет} (b4.north);
\draw[flow] (b5.west) -- (20,44) -- node[lbl,left,pos=0.35] {да} (b6.north);
\draw[flow] (b5.east) -- (60,44) -- node[lbl,right,pos=0.35] {нет} (b7.north);
% слияние линий в смещённых точках (ГОСТ 19.701-90, п. 4.2.3)
\draw (b3.south) |- (0,77);
\draw[flow] (b4.south) |- (0,81);
\draw[flow] (b6.south) |- (0,85);
\draw[flow] (b7.south) |- (0,89);
\draw[flow] (0,77) -- (end.north);
\end{tikzpicture}

    \caption{Схема алгоритма вычисления функции $F$}
    \label{fig:scheme-f}
\end{figure}

\begin{figure}[H]
    \centering
    \begin{tikzpicture}[x=1mm,y=-1mm]
\node[terminator] (start) at (0,0) {$S(i,j,c)$};
\node[decision] (b9) at (0,24) {$c>9$ или $c>j$};
\node[process] (b10) at (-38,50) {$t=0$};
\node[predefined] (b11) at (38,50) {$a=F(i-1,j-c)$};
\node[predefined] (b12) at (38,68) {$b=S(i,j,c+1)$};
\node[process] (b13) at (38,86) {$t=a+b$};
\node[terminator] (end) at (0,116) {Возврат: $t$};
\draw[flow] (start) -- (b9);
\draw[flow] (b9.west) -- node[lbl,above] {да} (-38,24) -- (b10.north);
\draw[flow] (b9.east) -- node[lbl,above] {нет} (38,24) -- (b11.north);
\draw[flow] (b11) -- (b12);
\draw[flow] (b12) -- (b13);
% слияние линий в смещённых точках (ГОСТ 19.701-90, п. 4.2.3)
\draw (b10.south) |- (0,97);
\draw[flow] (b13.south) |- (0,102);
\draw[flow] (0,97) -- (end.north);
\end{tikzpicture}

    \caption{Схема алгоритма вычисления функции $S$}
    \label{fig:scheme-s}
\end{figure}

\bigskip
\noindent\begin{minipage}{\dimexpr\linewidth-0.5mm\relax}
\begin{lstlisting}[language=Python,caption={Псевдокод рекурсивного алгоритма},label={lst:rec}]
def F(i, j):
    if i == 0:
        if j == 0:
            r = 1
        else:
            r = 0
    elif j == 0:
        r = 0
    else:
        r = S(i, j, 0)
    return r
def S(i, j, c):
    if c > 9 or c > j:
        t = 0
    else:
        a = F(i - 1, j - c)
        b = S(i, j, c + 1)
        t = a + b
    return t
\end{lstlisting}
\end{minipage}
\medskip

Вызов функции обозначен на схемах символом «Предопределённый процесс»
(строки 10, 16 и 17 листинга~\ref{lst:rec}). Выполнение вызова передаёт
аргументы вызываемой функции, а результат, возвращённый вызванной функцией
в строке 11 или 19, присваивается переменной $r$, $a$ или $b$. Повторные вызовы $F$
с одинаковыми аргументами вычисляются заново, поскольку сохранять результаты
в массиве запрещено.

В графовых моделях исследовательской части оператор обозначается номером
строки листинга~\ref{lst:iter} или~\ref{lst:rec}. Строки с заголовками
функций и строки \texttt{else} операторами не считаются.
